\documentclass[aspectratio=169]{beamer}

\usetheme{Madrid}
\usecolortheme{default}
\setbeamertemplate{navigation symbols}{}

\title[kira-bt]{KIRA-BT}
\subtitle{3x Faster Variant Annotation}
\author{Kira Bio Tools}
\date{}

\begin{document}

\begin{frame}
  \titlepage
  \vspace{0.5em}
  \begin{center}
    \small Remove the annotation bottleneck in genomics pipelines.
  \end{center}
\end{frame}

\begin{frame}{The Problem}
\begin{itemize}
  \item Variant annotation is the slow step in genomics pipelines.
  \item Industry standard: \texttt{bcftools annotate}.
  \item Performance is I/O-bound and scales poorly with volume.
  \item Slow annotation delays downstream analysis and reporting.
\end{itemize}
\end{frame}

\begin{frame}{The Solution}
\begin{itemize}
  \item \textbf{kira-bt} is a drop-in replacement for \texttt{bcftools annotate}.
  \item 3x faster variant annotation, no pipeline changes required.
  \item Same formats: VCF / BGZF / CSI / TBI.
  \item Benchmark: 427,000 variants/sec vs 80k--120k variants/sec.
\end{itemize}
\end{frame}

\begin{frame}{Technology Overview}
\begin{itemize}
  \item Prebuilt annotation index for fast lookups.
  \item Streaming input to avoid random I/O.
  \item Batch processing for throughput.
  \item Multithreaded execution across CPU cores.
\end{itemize}
\end{frame}

\begin{frame}{Benchmarks}
\begin{itemize}
  \item End-to-end measurement (read + annotate + write).
\end{itemize}
\vspace{0.5em}
\begin{tabular}{l r}
  \textbf{Tool} & \textbf{Speed} \\\hline
  bcftools annotate & 80k--120k variants/sec \\
  kira-bt & 427k variants/sec \\
\end{tabular}
\vspace{0.5em}
\begin{itemize}
  \item 3x speedup on real datasets.
\end{itemize}
\end{frame}

\begin{frame}{Time Savings}
\begin{itemize}
  \item Per whole genome (~5M variants):
  \begin{itemize}
    \item bcftools: 50s
    \item kira-bt: 12s
    \item 38s saved
  \end{itemize}
  \item 1,000 genomes: 10.6 hours saved
  \item 10,000 genomes: 4.4 days saved
\end{itemize}
\end{frame}

\begin{frame}{Cost Savings}
\begin{itemize}
  \item Lower compute cost per cohort.
  \item Higher throughput with existing hardware.
  \item Faster time-to-result improves operational efficiency.
\end{itemize}
\end{frame}

\begin{frame}{Market}
\begin{itemize}
  \item Clinical genomics labs
  \item Research institutes and universities
  \item Pharma and biotech companies
  \item National biobanks
\end{itemize}
\vspace{0.5em}
Growing data volumes, established pipelines, high demand for compatibility.
\end{frame}

\begin{frame}{Business Model}
\begin{itemize}
  \item Enterprise on-prem licenses
  \item Paid support and SLA
  \item Cloud images / managed runners
  \item OEM licensing to pipeline vendors
\end{itemize}
\end{frame}

\begin{frame}{Roadmap}
\begin{itemize}
  \item 6--12 months: performance improvements and stability.
  \item Full bcftools feature parity.
  \item Validation for clinical-grade usage.
\end{itemize}
\end{frame}

\begin{frame}{Team}
\begin{itemize}
  \item Founder-led, systems and bioinformatics background.
  \item Proven execution on performance-critical tooling.
  \item Focused on delivery and customer adoption.
\end{itemize}
\end{frame}

\begin{frame}{Ask}
\begin{itemize}
  \item Raising a pre-seed / seed round.
  \item Capital for productization and early customers.
  \item Value: remove the annotation bottleneck without pipeline changes.
\end{itemize}
\end{frame}

\end{document}
